\begin{table*}[htbp]
  \centering
  \caption{Component-level information pattern. The baseline cascade and the $L_1$ and MPC extensions consume no fault flag and no peer message; only the optional reshape supervisor steps outside this pattern.}
  \label{tab:component-pattern}
  \small
  \begin{tabular}{p{4.5cm}cccp{3.0cm}}
    \toprule
    Component & Detection? & Communication? & In main theorem? & Purpose \\
    \midrule
    Baseline cascade + identity \eqref{eq:Tff} & No & No & Yes & Passive fault absorption \\
    $L_1$ altitude augmentation & No & No & Preserves theorem & Parameter mismatch \\
    MPC tension ceiling & No & No & Preserves theorem & Tension constraint \\
    Reshape supervisor & Yes (latch) & Yes ($\lceil\log_2 N\rceil$ bits) & Optional only & Slot reassignment \\
    \bottomrule
  \end{tabular}
\end{table*}

We present the decentralized fault-tolerant controller for \cref{def:problem}. Each drone runs a three-layer cascade: an outer-loop PD with anti-swing slot shift, a single-step convex QP on the tilt-and-thrust envelope, and an inner-loop attitude PD. The central design element is $T_i^{\mathrm{ff}} = T_i$, encoding passive fault tolerance with no detection, communication, or reconfiguration. The reshape supervisor (\cref{sec:reshape-supervisor}) adds a latch detector and a per-fault broadcast; the optional $L_1$ augmentation (\cref{sec:extensions}) injects an adaptive altitude correction between the outer PD and the QP.

\begin{figure*}[t]
\centering
\begin{tikzpicture}[
    >={Stealth[length=2mm,width=1.6mm]},
    blk/.style={draw, thick, rounded corners=1pt, fill=blue!8,
                minimum height=11mm, minimum width=22mm,
                align=center, font=\footnotesize, inner sep=2pt},
    extblk/.style={draw, thick, rounded corners=1pt, fill=orange!18,
                   minimum height=10mm, minimum width=26mm,
                   align=center, font=\footnotesize, inner sep=2pt},
    sens/.style={draw, thick, rounded corners=1pt, fill=green!10,
                 minimum height=9mm, minimum width=30mm,
                 align=center, font=\footnotesize, inner sep=2pt},
    plant/.style={draw, thick, rounded corners=1pt, fill=gray!12,
                  minimum height=11mm, minimum width=24mm,
                  align=center, font=\footnotesize, inner sep=2pt},
    sum/.style={draw, thick, circle, inner sep=0.5pt,
                minimum size=4.5mm, font=\scriptsize},
    sig/.style={font=\scriptsize, inner sep=1.2pt, fill=white},
    boxed/.style={draw, dashed, rounded corners=2pt, inner sep=4pt},
    flow/.style={->, thick},
    fftap/.style={->, thick, red!70!black},
    senseflow/.style={->, thick, blue!55!black},
    extflow/.style={->, thick, dashed, orange!75!black}
  ]

  % --- Baseline cascade row (y = 0) ---
  \node[sig, font=\scriptsize] (refin) at (-8.6, 0)
        {$\mathbf{p}_L^d,\,\mathbf{v}_L^d$};
  \node[blk] (pd)    at (-6.0, 0) {Outer PD\\$+$ anti-swing\\slot shift};
  \node[sum] (sumL1) at (-3.0, 0) {$+$};
  \node[blk] (qp)    at ( 0.0, 0) {Single-step QP\\tilt + thrust\\envelope};
  \node[blk] (att)   at ( 3.4, 0) {Attitude PD\\(ZYX Euler)};
  \node[plant] (plant) at ( 6.7, 0) {Drone $i$ + cable\\(reduced model)};

  % horizontal cascade flow
  \draw[flow] (refin) -- (pd);
  \draw[flow] (pd) -- node[sig, above] {$\mathbf{a}_{\mathrm{target}}$}
        (sumL1);
  \draw[flow] (sumL1) -- (qp);
  \draw[flow] (qp) -- node[sig, above] {$\mathbf{a}_d^\star$} (att);
  \draw[flow] (att) -- node[sig, above] {$f_i,\,\boldsymbol{\tau}_i$}
        (plant);

  % --- Feed-forward identity tap (sumT directly above QP) ---
  \node[sum, fill=red!10] (sumT) at (0, 1.7) {$+$};
  % Identity label: placed ABOVE sumT (out of any tap path);
  % no white fill so it does not look like a small box.
  \node[anchor=south, font=\scriptsize\itshape, red!70!black,
        inner sep=2pt] at (sumT.north)
        {$T_i^{\mathrm{ff}}\!=\!T_i$\,(identity)};
  \draw[fftap] (sumT.south) -- (qp.north);

  % --- Sensor band (y = -2.6, centre at x = 1.7) ---
  % x = 1.7 sits in the clear gap between QP (right edge x=1.1)
  % and att (left edge x=2.3), so sense.north runs vertically up
  % to sumT without crossing any block.
  \node[sens] (sense) at (1.7, -2.6)
        {Local sensors\\$x_i,\;\; T_i,\;\; v_L$};

  % Plant -> sensors: |- (vertical down then horizontal left)
  \draw[senseflow] (plant.south) |- (sense.east);

  % Sensors -> PD: horizontal left at y = -2.6, then vertical up
  % at x = -6.0.  Single draw call -> ONE arrowhead at pd.south.
  \draw[senseflow] (sense.west)
        -- node[sig, pos=0.65, above] {$x_i,\,v_L$}
        (-6.0, -2.6) -- (pd.south);

  % T_i tap (red) — sense.NORTH straight up, then horizontal
  % left to sumT.east.  Single draw call -> ONE arrowhead at
  % sumT.east.  Vertical at x=1.7 clears QP (right edge x=1.1)
  % and att (left edge x=2.3); horizontal at y=1.7 sits well
  % above all baseline blocks (top edges at y ~= 0.55).
  % T_i label placed near the start of the vertical (well below
  % the cascade row), with no fill so it does not look like a
  % small block sitting on the cascade.
  \draw[fftap] (sense.north)
        -- node[anchor=west, font=\scriptsize\itshape,
                red!70!black, pos=0.20, inner sep=1pt]
                {$T_i$ (own tension)}
        (sense.north |- sumT.east) -- (sumT.east);

  % --- Bracket the baseline ---
  \begin{scope}[on background layer]
    \node[boxed,
          fit=(refin) (pd) (sumL1) (qp) (att) (plant) (sumT) (sense)]
         (baseline) {};
  \end{scope}
  \node[font=\scriptsize\itshape, anchor=south west]
        at (baseline.north west)
        {Baseline cascade --- continuous time, all signals local};

  % --- Extension band (y = -4.6) ---
  \node[extblk] (reshape) at (-6.0, -4.6)
        {Reshape supervisor\\$\lceil\log_2 N\rceil$-bit broadcast};
  \node[extblk] (l1) at (-3.0, -4.6)
        {$L_1$ adaptive\\altitude augmentation};
  \node[extblk] (mpc) at ( 0.0, -4.6)
        {Receding-horizon MPC\\(replaces single-step QP)};

  % All three extension injection arrows are strictly vertical.
  \draw[extflow] (reshape.north)
        -- node[sig, pos=0.78, right=2pt]
              {$\boldsymbol{\delta}_i^{\mathrm{ovr}}$}
        (pd.south);
  \draw[extflow] (l1.north)
        -- node[sig, pos=0.78, right=2pt] {$u_{\mathrm{ad}}$}
        (sumL1.south);
  \draw[extflow] (mpc.north)
        -- node[sig, pos=0.78, right=2pt] {replaces QP}
        (qp.south);

  % --- Bracket the extension layers ---
  \begin{scope}[on background layer]
    \node[boxed, fit=(reshape) (l1) (mpc), draw=orange!70!black]
         (ext) {};
  \end{scope}
  \node[font=\scriptsize\itshape, orange!70!black,
        anchor=south west]
        at (ext.north west)
        {Extension layers (\cref{sec:extensions})};

\end{tikzpicture}
\caption{Per-drone control cascade. The baseline (blue, framed top) is
a three-stage continuous-time pipeline: outer-loop PD with anti-swing
slot shift, a single-step convex QP that projects
$\mathbf{a}_{\mathrm{target}}$ onto the tilt-and-thrust envelope, and an
Euler-angle attitude PD. The non-standard element is the red
feed-forward tap: drone $i$ routes its own measured tension $T_i$
directly into the altitude thrust feed-forward,
$T_i^{\mathrm{ff}} = T_i$. This single identity is the carrier of
hybrid practical-ISS across cable-severance events
(\cref{thm:hybrid-stability}); the surrounding cascade is otherwise
textbook. Three optional layers (orange, framed bottom) hook in at
three precisely identified injection points: the $L_1$ adaptive
correction $u_{\mathrm{ad}}$ enters the altitude channel between the
outer PD and the QP; the receding-horizon MPC replaces the single-step
QP entirely; the reshape supervisor overrides the formation slot
$\boldsymbol{\delta}_i$ on receipt of a single
$\lceil\log_2 N\rceil$-bit broadcast.}
\label{fig:controller-cascade}
\end{figure*}


Each drone is pre-assigned a hovering slot displaced vertically above the payload reference point and laterally by a fixed formation offset,
\begin{equation}
  \vect p_{\mathrm{slot},i}(t)
    \;=\; \vect p_L^d(t) + \boldsymbol\delta_i
        + d_{\mathrm{rope}}\,\hat{\vect e}_3,
  \label{eq:slot-pos}
\end{equation}
where $\boldsymbol\delta_i \in \R^3$ is the formation offset (overridable at run time), $d_{\mathrm{rope}} = \SI{1.25}{\metre}$ is the vertical slot elevation matching $L$, and $\hat{\vect e}_3$ is the upward unit vector. The slot velocity is $\vect v_{\mathrm{slot},i}(t) = \vect v_L^d(t)$. To damp the payload pendulum passively, we shift the slot relative to the payload's horizontal velocity. With $\vect v_{L,\perp} \triangleq [v_{L,x}, v_{L,y}, 0]^{\top}$, the horizontal swing correction is
\begin{equation}
  \vect s_i(t) =
    \begin{cases}
      -k_{\mathrm{swing}}\,\vect v_{L,\perp}
        & \text{if }
          k_{\mathrm{swing}}\|\vect v_{L,\perp}\| \le s_{\max}, \\[2pt]
      -s_{\max}\,\hat{\vect v}_{L,\perp}
        & \text{otherwise,}
    \end{cases}
  \label{eq:swing-shift}
\end{equation}
where $\hat{\vect v}_{L,\perp}$ is the unit vector of $\vect v_{L,\perp}$. The shifted slot position is $\vect p_{\mathrm{slot},i}^{\mathrm{dyn}}(t) = \vect p_{\mathrm{slot},i}(t) + \vect s_i(t)$. therefore, the position and velocity errors are
\begin{align}
    \label{eq:slot-error}
  \vect e_p(t) &= \vect p_{\mathrm{slot},i}^{\mathrm{dyn}}(t)
                 - \vect p_i(t),\\
  \vect e_v(t) &= \vect v_{\mathrm{slot},i}(t) - \vect v_i(t).\nonumber
\end{align}

The outer loop computes a target body acceleration,
\begin{align}
\label{eq:a-target}
  \vect a_{\mathrm{target}}(t)
    =& \underbrace{%
        \begin{bmatrix}
          K_p^{xy}\, e_{p,x} + K_d^{xy}\, e_{v,x} \\[2pt]
          K_p^{xy}\, e_{p,y} + K_d^{xy}\, e_{v,y} \\[2pt]
          K_p^{z}\;\; e_{p,z} + K_d^{z}\;\; e_{v,z}
        \end{bmatrix}}_{\vect a_{\mathrm{track}}}\\
    &+\; w_{\mathrm{swing}}
      \begin{bmatrix}
        -k_{\mathrm{swing}}\, v_{L,x} \\[2pt]
        -k_{\mathrm{swing}}\, v_{L,y} \\[2pt]
        0
      \end{bmatrix}.\nonumber
\end{align}
The altitude gain $K_p^z = 100$ exceeds the horizontal gain $K_p^{xy} = 30$ because vertical rope forces dominate the disturbance budget; the formal pole and Lyapunov-matrix analysis is in \cref{lem:altitude-iss}. When $L_1$ is active, its altitude correction $u_{\mathrm{ad}}(t)$ is added to the $z$-component of $\vect a_{\mathrm{target}}$ between the outer PD and the QP (\cref{sec:extensions}); horizontal channels are unmodified.

The target acceleration $\vect a_{\mathrm{target}}$ may lie outside the tilt-and-thrust actuator envelope. We project onto the feasible set with a single-step convex QP that also regularises large commands:
\begin{subequations}
\label{eq:qp}
\begin{align}
  \vect a_d^\star
    &= \operatorname*{argmin}_{\vect a_d \in \R^3}
      \; w_t\,\bigl\|\vect a_d - \vect a_{\mathrm{target}}\bigr\|^2
      + w_e\,\|\vect a_d\|^2,
    \label{eq:qp-cost} \\
  \text{s.t.}\;\;
    &\abs{a_{d,x}},\;\abs{a_{d,y}}
      \le a_{\mathrm{tilt}} \triangleq g\tan\theta_{\max},
    \label{eq:qp-tilt} \\
    &a_{z,\min}(t) \le a_{d,z} \le a_{z,\max}(t),
    \label{eq:qp-vert}
\end{align}
\end{subequations}
where the vertical bounds incorporate the rope tension feed-forward,
\begin{align}
\label{eq:az-bounds}
  a_{z,\min/\max}(t) &= \frac{f_{\min\max} - T_i^{\mathrm{ff}}(t)}{m_i} - g.
\end{align}
The vertical interval in \eqref{eq:qp-vert} is non-empty whenever $f_{\min} < f_{\max}$; hover feasibility holds in all scenarios, and Doctrine~D1 is satisfied at the $0.10\%$ worst QP-transition fraction (\cref{tab:domain-audit}). For the canonical campaigns, the feed-forward is the identity from $t=0$:
\begin{equation}
  \boxed{T_i^{\mathrm{ff}}(t) \;=\; T_i(t).}
  \label{eq:Tff}
\end{equation}
A smoothstep pickup ramp $\alpha(\tau) = 3\tau^2 - 2\tau^3$ handles cold-start engagement; it is bypassed in every reported campaign.

Given the QP solution $a_{d,z}^\star$, the commanded thrust is
\begin{equation}
\label{eq:thrust}
  f_i \;=\; \operatorname{clamp}\!\bigl(\,
      m_i(g + a_{d,z}^\star) + T_i^{\mathrm{ff}},\;
      f_{\min},\; f_{\max}\bigr).
\end{equation}
Under the small-angle approximation, the horizontal acceleration commands map to desired pitch and roll angles,
\begin{align}
\label{eq:attitude-des}
  \theta_{\mathrm{pitch}}^d
    &= \operatorname{sat}\!\!\left(\frac{a_{d,x}^\star}{g},\;\pm\theta_{\max}\right),\\
  \theta_{\mathrm{roll}}^d
    &= \operatorname{sat}\!\!\left(-\frac{a_{d,y}^\star}{g},\;\pm\theta_{\max}\right).\nonumber
\end{align}
With body angular velocity $\vect\omega_i^{\mathcal{B}} = \mat R_i^{\top}\, \vect\omega_i^{\mathcal{W}}$, the body torques are
\begin{align}
  \tau_{i,x}
    &= K_p^{\mathrm{att}}\!\left(
         \theta_{\mathrm{roll}}^d - \phi_{\mathrm{roll}}
       \right)
      - K_d^{\mathrm{att}}\,\omega_{i,x}^{\mathcal{B}},
    \label{eq:att-pd-roll} \\[2pt]
  \tau_{i,y}
    &= K_p^{\mathrm{att}}\!\left(
         \theta_{\mathrm{pitch}}^d - \theta_{\mathrm{pitch}}
       \right)
      - K_d^{\mathrm{att}}\,\omega_{i,y}^{\mathcal{B}},
    \label{eq:att-pd-pitch} \\[2pt]
  \tau_{i,z}
    &= -\frac{K_p^{\mathrm{att}}}{2}\,(R_{21} - R_{12})
          - K_d^{\mathrm{att}}\,\omega_{i,z}^{\mathcal{B}},
    \label{eq:att-pd-yaw}
\end{align}
each saturated to $[-\tau_{\max}, \tau_{\max}]$. The law \eqref{eq:att-pd-roll}--\eqref{eq:att-pd-yaw} uses Euler-angle error directly rather than a geodesic $\SOthree$ map; this is valid in the small-tilt regime enforced by \eqref{eq:qp-tilt}. The output wrench is
\begin{equation}
  \vect F_i^{\mathcal{W}} = f_i\,\mat R_i\,\hat{\vect e}_3,
  \qquad
  \vect\tau_i^{\mathcal{W}} = \mat R_i\,
    [\tau_{i,x},\; \tau_{i,y},\; \tau_{i,z}]^{\top},
  \label{eq:wrench}
\end{equation}
applied at the same continuous-time evaluation instant as the upstream computations. \Cref{tab:baseline-params} collects the canonical controller parameters used throughout all simulation campaigns.

\begin{table}[t]
  \centering
  \caption{Proposed controller parameters (canonical values used
  in all simulation campaigns).}
  \label{tab:baseline-params}
  \small
  \renewcommand{\arraystretch}{1.05}
  \begin{tabular}{lrl}
    \toprule
    Parameter & Value & Role \\
    \midrule
    $K_p^{xy},\; K_d^{xy}$
      & $30,\;15$
      & horizontal outer loop \\
    $K_p^{z},\;\; K_d^{z}$
      & $100,\;24$
      & altitude outer loop \\
    $K_p^{\mathrm{att}},\; K_d^{\mathrm{att}}$
      & $25,\;4$
      & attitude inner loop \\
    $k_{\mathrm{swing}}$
      & $0.8$
      & anti-swing gain \\
    $w_{\mathrm{swing}}$
      & $0.3$
      & anti-swing weight in $\vect a_{\mathrm{target}}$ \\
    $s_{\max}$
      & \SI{0.3}{\metre}
      & slot-shift saturation \\
    $d_{\mathrm{rope}}$
      & \SI{1.25}{\metre}
      & slot $z$-offset (nominal rope length) \\
    $w_t,\; w_e$
      & $1.0,\;0.02$
      & QP tracking / effort weights \\
    $\theta_{\max}$
      & \SI{0.6}{\radian}
      & tilt ceiling (${\approx}34^{\circ}$) \\
    $f_{\min},\; f_{\max}$
      & $0,\;\SI{150}{\newton}$
      & thrust envelope \\
    $\tau_{\max}$
      & \SI{10}{\newton\metre}
      & torque saturation \\
    $T_{\mathrm{thresh}}$
      & \SI{0.3}{\newton}
      & pickup-detection threshold \\
    $\bar{T}$
      & \SI{19.62}{\newton}
      & nominal per-drone tension ($m_L g / N$) \\
    $K_{p,T}$
      & $1.5$
      & pickup-phase tension correction gain \\
    $T_{\mathrm{ramp}}$
      & \SI{1.0}{\second}
      & pickup smoothstep window \\
    $\Delta t_{L_1}$
      & \SI{0.2}{\milli\second}
      & $L_1$ discrete update period \\
    \bottomrule
  \end{tabular}
\end{table}